% =================================================================
%  CSE 815 -- Machine Learning
%  Module 1, Week 3, Part 2
%  Overfitting, and Regularization for Linear and Logistic
%  Regression
% =================================================================
\documentclass[aspectratio=169,10pt]{beamer}
\usepackage{cse815}

\title[CSE 815 -- Week 3, Part 2]{Overfitting and Regularization}
\subtitle{Why a Perfect Fit Is a Warning Sign, and What to Do About It}
\author{Rokan Uddin Faruqui\\ \small Professor, Department of CSE}
\institute{\small University of Chittagong \\ \texttt{rufaruqui@cu.ac.bd}}
\date[Week 3 / Part 2]{Week 3 $\cdot$ Part 2}

\begin{document}

\begin{frame}[plain,noframenumbering]
  \titlepage
\end{frame}

% ---------------------------------------------------------------
\begin{frame}{The question Week 2 left open}
  \begin{columns}[T,onlytextwidth]
    \begin{column}{0.5\textwidth}
      \begin{block}{Where we stand}
        We can now fit
        \[ f_{\vw,b}(\vx)=\vw\cdot\vx+b \quad\text{(a number)} \]
        \[ f_{\vw,b}(\vx)=g(\vw\cdot\vx+b) \quad\text{(a class)} \]
        and we can manufacture as many features as we like --- $x^2$, $x^3$, $x_1x_2$, \dots
      \end{block}
      \footnotesize
      Week 2 ended with a warning: adding features always lowers the training cost. So the
      training cost cannot tell us when to stop. Today we look at what actually goes wrong,
      and at the standard remedy.
    \end{column}
    \begin{column}{0.46\textwidth}
      \begin{keyidea}[Today]
        \footnotesize
        \begin{enumerate}
          \item Underfitting, overfitting, generalization 
          \item Overfitting in classification 
          \item Three ways to address it 
          \item The regularized cost and $\lambda$  
          \item Regularized linear regression  
          \item Regularized logistic regression  
        \end{enumerate}
      \end{keyidea}
      \footnotesize This closes Module 1: everything here applies unchanged to the neural
      networks of Module 2.
    \end{column}
  \end{columns}
\end{frame}

% ===============================================================
\section{The Problem of Overfitting}

\begin{frame}{Three fits to the same seven flats}
  \begin{columns}[T,onlytextwidth]
    \begin{column}{0.33\textwidth}
      \centering {\scriptsize \textbf{Underfit} --- $f = w_1x+b$}\par
      \begin{tikzpicture}
        \begin{axis}[width=4.3cm,height=3.5cm,xmin=0.4,xmax=3.0,ymin=10,ymax=70,
          xlabel={size},ylabel={price},grid=major,grid style={culight},
          tick label style={font=\tiny},label style={font=\tiny}]
          \addplot[only marks,mark=*,mark size=1.4pt,cublue,forget plot]
            coordinates {(0.6,20)(0.9,26)(1.2,36)(1.5,38)(1.9,49)(2.3,50)(2.8,58)};
          \addplot[curose,very thick,domain=0.5:2.9,forget plot]{12.2236+17.0924*x};
        \end{axis}
      \end{tikzpicture}\par
      {\tiny $J_{\text{train}} = 3.19$}
    \end{column}
    \begin{column}{0.33\textwidth}
      \centering {\scriptsize \textbf{Good fit} --- add $x^2$}\par
      \begin{tikzpicture}
        \begin{axis}[width=4.3cm,height=3.5cm,xmin=0.4,xmax=3.0,ymin=10,ymax=70,
          xlabel={size},ylabel={price},grid=major,grid style={culight},
          tick label style={font=\tiny},label style={font=\tiny}]
          \addplot[only marks,mark=*,mark size=1.4pt,cublue,forget plot]
            coordinates {(0.6,20)(0.9,26)(1.2,36)(1.5,38)(1.9,49)(2.3,50)(2.8,58)};
          \addplot[cuteal,very thick,domain=0.5:2.9,samples=60,forget plot]
            {2.8677+30.6696*x-4.0080*x^2};
        \end{axis}
      \end{tikzpicture}\par
      {\tiny $J_{\text{train}} = 1.44$}
    \end{column}
    \begin{column}{0.33\textwidth}
      \centering {\scriptsize \textbf{Overfit} --- up to $x^6$}\par
      \begin{tikzpicture}
        \begin{axis}[width=4.3cm,height=3.5cm,xmin=0.4,xmax=3.0,ymin=10,ymax=70,
          xlabel={size},ylabel={price},grid=major,grid style={culight},
          tick label style={font=\tiny},label style={font=\tiny}]
          \addplot[only marks,mark=*,mark size=1.4pt,cublue,forget plot]
            coordinates {(0.6,20)(0.9,26)(1.2,36)(1.5,38)(1.9,49)(2.3,50)(2.8,58)};
          \addplot[cuamber,very thick,domain=0.58:2.82,samples=200,forget plot]
            {790.0825-3765.3895*x+7153.0487*x^2-6792.7813*x^3+3441.0759*x^4-886.2694*x^5+91.0346*x^6};
        \end{axis}
      \end{tikzpicture}\par
      {\tiny $J_{\text{train}} = 0.00$}
    \end{column}
  \end{columns}
  \vspace{1mm}
  \footnotesize
  Seven points, seven parameters: the degree-$6$ curve passes through every one, so the
  training cost is exactly zero --- and it predicts that a $2.6$k~sq~ft flat is worth
  \emph{less} than a $2.1$k~sq~ft one.
  \begin{keyidea}
    The lowest training cost is the worst model. What we want is not a fit but
    \textbf{generalization}: low cost on examples the model has never seen.
  \end{keyidea}
\end{frame}

\begin{frame}{The vocabulary}
  \footnotesize
  \begin{center}
  \renewcommand{\arraystretch}{1.25}
  \begin{tabular}{@{}L{2.3cm}L{4.3cm}L{4.6cm}@{}}
    \toprule
    & \textbf{Underfitting} & \textbf{Overfitting}\\
    \midrule
    Also called      & high bias & high variance\\
    Symptom          & poor on training \emph{and} on new data & excellent on training,
                       poor on new data\\
    What it means    & a strong preconception the data contradicts &
                       the noise has been fitted as if it were signal\\
    Cause            & too few features, too simple a form & too many features for the
                       amount of data\\
    \bottomrule
  \end{tabular}
  \end{center}
  \begin{columns}[T,onlytextwidth]
    \begin{column}{0.5\textwidth}
      \begin{defnbox}[Why ``variance'']
        \footnotesize Re-collect the training data and refit. The degree-$6$ curve would come
        out completely different every time; the straight line would barely move. High
        variance means the fitted function depends violently on which examples you drew.
      \end{defnbox}
    \end{column}
    \begin{column}{0.46\textwidth}
      \begin{pitfall}
        \footnotesize
        You cannot detect overfitting from $J_{\text{train}}$ --- by construction it looks
        best when the problem is worst. You need data the model did not train on. Module 2
        builds this into a train / cross-validation / test procedure.
      \end{pitfall}
    \end{column}
  \end{columns}
\end{frame}

\begin{frame}{The same thing in classification}
  \begin{columns}[T,onlytextwidth]
    \begin{column}{0.33\textwidth}
      \centering {\scriptsize \textbf{Underfit}: $z=w_1x_1+w_2x_2+b$}\par
      \begin{tikzpicture}
        \begin{axis}[width=4.3cm,height=3.6cm,xmin=0,xmax=4,ymin=0,ymax=4,
          xlabel={$x_1$},ylabel={$x_2$},axis on top,
          tick label style={font=\tiny},label style={font=\tiny}]
          \addplot[only marks,mark=o,mark size=1.6pt,cublue,thick,forget plot]
            coordinates {(0.5,0.9)(0.9,1.4)(1.4,0.8)(0.6,2.0)(2.2,0.7)(1.1,2.2)(1.5,2.4)(2.0,1.75)};
          \addplot[only marks,mark=x,mark size=2pt,curose,thick,forget plot]
            coordinates {(2.6,2.6)(3.2,1.5)(1.8,3.4)(3.5,3.0)(2.9,1.2)(1.3,3.6)(2.4,1.0)};
          \addplot[cublue,very thick,domain=0.05:3.25,forget plot]{3.9-1.2*x};
        \end{axis}
      \end{tikzpicture}\par
      {\tiny 3 training points wrong}
    \end{column}
    \begin{column}{0.33\textwidth}
      \centering {\scriptsize \textbf{Good fit}: add $x_1^2,x_2^2$}\par
      \begin{tikzpicture}
        \begin{axis}[width=4.3cm,height=3.6cm,xmin=0,xmax=4,ymin=0,ymax=4,
          xlabel={$x_1$},ylabel={$x_2$},axis on top,
          tick label style={font=\tiny},label style={font=\tiny}]
          \addplot[only marks,mark=o,mark size=1.6pt,cublue,thick,forget plot]
            coordinates {(0.5,0.9)(0.9,1.4)(1.4,0.8)(0.6,2.0)(2.2,0.7)(1.1,2.2)(1.5,2.4)(2.0,1.75)};
          \addplot[only marks,mark=x,mark size=2pt,curose,thick,forget plot]
            coordinates {(2.6,2.6)(3.2,1.5)(1.8,3.4)(3.5,3.0)(2.9,1.2)(1.3,3.6)(2.4,1.0)};
          \addplot[cuteal,very thick,domain=0.05:3.15,samples=80,forget plot]
            {3.8-0.30*x-0.30*x^2};
        \end{axis}
      \end{tikzpicture}\par
      {\tiny 1 training point wrong}
    \end{column}
    \begin{column}{0.33\textwidth}
      \centering {\scriptsize \textbf{Overfit}: high-order terms}\par
      \begin{tikzpicture}
        \begin{axis}[width=4.3cm,height=3.6cm,xmin=0,xmax=4,ymin=0,ymax=4,
          xlabel={$x_1$},ylabel={$x_2$},axis on top,
          tick label style={font=\tiny},label style={font=\tiny}]
          \addplot[only marks,mark=o,mark size=1.6pt,cublue,thick,forget plot]
            coordinates {(0.5,0.9)(0.9,1.4)(1.4,0.8)(0.6,2.0)(2.2,0.7)(1.1,2.2)(1.5,2.4)(2.0,1.75)};
          \addplot[only marks,mark=x,mark size=2pt,curose,thick,forget plot]
            coordinates {(2.6,2.6)(3.2,1.5)(1.8,3.4)(3.5,3.0)(2.9,1.2)(1.3,3.6)(2.4,1.0)};
          \addplot[cuamber,very thick,domain=0.62:3.3,samples=200,forget plot]
            {3.8-0.30*x-0.30*x^2+0.6*sin(deg(4.5*x))};
        \end{axis}
      \end{tikzpicture}\par
      {\tiny 0 training points wrong}
    \end{column}
  \end{columns}
  \vspace{1mm}
  \footnotesize
  Same fifteen patients, three boundaries. Judged on the training set the right-hand model is
  the best of the three --- and it is the one you would least want in a clinic. It has bent
  itself around individual examples, so a new patient landing in one of its inlets is
  classified by an accident of who happened to be in the training set.
  \begin{keyidea}
    Overfitting is not about regression or classification. It is about a model flexible enough
    to memorise the examples instead of describing the pattern behind them.
  \end{keyidea}
\end{frame}

% ===============================================================
\section{Addressing Overfitting}

\begin{frame}{Three remedies}
  \footnotesize
  \begin{columns}[T,onlytextwidth]
    \begin{column}{0.32\textwidth}
      \begin{block}{1. Collect more data}
        The most effective fix when it is available. With enough examples even a
        high-order polynomial cannot wiggle --- there is a point in the way of every
        detour.
        \par\vspace{1mm}
        \textbf{Limit:} often impossible. There may only be $200$ patients.
      \end{block}
    \end{column}
    \begin{column}{0.32\textwidth}
      \begin{block}{2. Use fewer features}
        Select a subset --- by domain knowledge, or systematically. Fewer parameters, less
        room to memorise.
        \par\vspace{1mm}
        \textbf{Limit:} you throw away information. If all $100$ features carry a little
        signal, discarding $90$ discards signal with the noise.
      \end{block}
    \end{column}
    \begin{column}{0.32\textwidth}
      \begin{block}{3. Regularization}
        Keep every feature, but force the weights to stay small. A gentler version of
        feature selection: the model may use a feature, but only if it earns its place.
        \par\vspace{1mm}
        \textbf{This is the subject of the rest of the lecture.}
      \end{block}
    \end{column}
  \end{columns}
  \vspace{1mm}
  \begin{keyidea}
    Why do small weights help? A large $w_j$ makes the prediction swing wildly for a small
    change in $x_j$ --- that is exactly the wiggle. In the degree-$6$ fit two slides back,
    $\sum_j w_j^2 \approx 1.2\times10^{8}$, with $|w_2| \approx 7153$. The wiggle and the
    huge coefficients are the same phenomenon.
  \end{keyidea}
\end{frame}

% ===============================================================
\section{The Cost Function with Regularization}

\begin{frame}{Penalising large weights}
  \small
  Add a term to the cost that grows when the weights do:
  \[
    J(\vw,b) \;=\; \underbrace{\frac{1}{2m}\sum_{i=1}^{m}\bigl(f_{\vw,b}(\xii{i})-\yii{i}\bigr)^2}
      _{\text{fit the data}}
    \;+\; \underbrace{\frac{\lambda}{2m}\sum_{j=1}^{n} w_j^{2}}_{\text{keep } \vw \text{ small}}
  \]
  \begin{columns}[T,onlytextwidth]
    \begin{column}{0.5\textwidth}
      \footnotesize
      \begin{itemize}
        \item $\lambda \ge 0$ is the \textbf{regularization parameter}. You choose it; it is
              not learned.
        \item Dividing by the same $2m$ makes a good $\lambda$ less sensitive to the size of
              the training set.
        \item The sum runs $j=1,\dots,n$: \textbf{$b$ is not penalised}. Shrinking the
              intercept towards zero says the output should be near zero by default, which is
              a claim about the world, not about complexity. In practice it makes almost no
              difference; convention is to leave it out.
      \end{itemize}
    \end{column}
    \begin{column}{0.46\textwidth}
      \begin{defnbox}[The trade-off]
        \footnotesize
        Minimising this cost balances two goals that pull in opposite directions:
        \begin{itemize}
          \item fit the training data,
          \item keep every $w_j$ small.
        \end{itemize}
        $\lambda$ sets the exchange rate between them.
      \end{defnbox}
      \footnotesize
      This penalty is called \textbf{L2} regularization, after the $\ell_2$ norm
      $\|\vw\|_2^2=\sum_j w_j^2$. In statistics the resulting model is \emph{ridge
      regression}.
    \end{column}
  \end{columns}
\end{frame}

\begin{frame}{What $\lambda$ does, at the two extremes}
  \begin{columns}[T,onlytextwidth]
    \begin{column}{0.5\textwidth}
      \footnotesize
      \begin{block}{$\lambda = 0$}
        The penalty vanishes; we are back to the unregularized fit. The degree-$6$ curve
        returns, with $J_{\text{train}}=0$ and coefficients in the thousands.
        \hfill\textbf{Overfits.}
      \end{block}
      \begin{block}{$\lambda \to \infty$ (say $10^{10}$)}
        The only way to keep $J$ small is $w_j \approx 0$ for every $j$, leaving
        $f_{\vw,b}(\vx) \approx b$. The model predicts the same number for every input --- a
        horizontal line. \hfill\textbf{Underfits.}
      \end{block}
      \begin{block}{In between}
        Some intermediate $\lambda$ balances the two. Finding it is a search, like the search
        for $\alpha$; Module 2 gives the principled procedure.
      \end{block}
    \end{column}
    \begin{column}{0.48\textwidth}
      \centering {\scriptsize The same degree-$6$ model, with and without a penalty}\par
      \begin{tikzpicture}
        \begin{axis}[width=6.2cm,height=4.3cm,xmin=0.4,xmax=3.0,ymin=10,ymax=70,
          xlabel={size (k sq ft)},ylabel={price (lakh)},grid=major,grid style={culight},
          tick label style={font=\scriptsize},label style={font=\scriptsize},
          legend style={font=\tiny,at={(0.97,0.03)},anchor=south east,draw=none,fill=white}]
          \addplot[only marks,mark=*,mark size=1.5pt,cublue,forget plot]
            coordinates {(0.6,20)(0.9,26)(1.2,36)(1.5,38)(1.9,49)(2.3,50)(2.8,58)};
          \addplot[cuamber,very thick,domain=0.58:2.82,samples=200]
            {790.0825-3765.3895*x+7153.0487*x^2-6792.7813*x^3+3441.0759*x^4-886.2694*x^5+91.0346*x^6};
          \addlegendentry{$\lambda=0$}
          \addplot[cuteal,very thick,domain=0.5:2.95,samples=120]
            {19.9618+2.0628*x+3.3158*x^2+3.3678*x^3+1.4930*x^4-2.0563*x^5+0.4026*x^6};
          \addlegendentry{$\lambda=1$}
        \end{axis}
      \end{tikzpicture}
    \end{column}
  \end{columns}
  \vspace{0.5mm}
  \footnotesize
  Both curves use \emph{all six} powers of $x$. With $\lambda=1$ the fitted weights collapse
  from $\sum_j w_j^2 \approx 1.2\times10^{8}$ to $33.2$, the largest from $7153$ to $3.37$,
  and $J_{\text{train}}$ rises from $0$ to $2.67$. The features were never removed --- they
  were merely made expensive.
\end{frame}

% ===============================================================
\section{Regularized Linear Regression}

\begin{frame}{The gradient, and the update}
  \footnotesize
  Differentiating the regularized cost adds one term, and only to the weights:
  \[
    \frac{\partial J}{\partial w_j}
      = \frac{1}{m}\sum_{i=1}^{m}\bigl(f_{\vw,b}(\xii{i})-\yii{i}\bigr)\xjj{i}{j}
        \;+\; \frac{\lambda}{m}w_j ,
    \qquad
    \frac{\partial J}{\partial b}
      = \frac{1}{m}\sum_{i=1}^{m}\bigl(f_{\vw,b}(\xii{i})-\yii{i}\bigr) .
  \]
  \begin{defnbox}[Regularized gradient descent]
    Repeat, simultaneously for $j=1,\dots,n$:
    \[
      w_j := w_j - \alpha\left[\frac{1}{m}\sum_{i=1}^{m}
        \bigl(f_{\vw,b}(\xii{i})-\yii{i}\bigr)\xjj{i}{j} + \frac{\lambda}{m}w_j\right],
      \qquad
      b := b - \alpha\,\frac{1}{m}\sum_{i=1}^{m}\bigl(f_{\vw,b}(\xii{i})-\yii{i}\bigr)
    \]
  \end{defnbox}
  \begin{pitfall}
    The $b$ update is unchanged. Adding $\tfrac{\lambda}{m}b$ to it is the single most common
    implementation error here, and it is silent --- the model still trains, just towards a
    slightly wrong intercept.
  \end{pitfall}
\end{frame}

\begin{frame}{What the extra term actually does}
  \footnotesize
  Collect the $w_j$ terms in the update:
  \[
    w_j := \underbrace{\left(1 - \alpha\frac{\lambda}{m}\right)}_{\text{shrink}} w_j
           \;-\; \underbrace{\alpha\frac{1}{m}\sum_{i=1}^{m}
             \bigl(f_{\vw,b}(\xii{i})-\yii{i}\bigr)\xjj{i}{j}}_{\text{the usual step}}
  \]
  \begin{columns}[T,onlytextwidth]
    \begin{column}{0.52\textwidth}
      Every iteration now does two things: it multiplies $w_j$ by a number slightly less than
      one, then takes the ordinary gradient step.
      \par\vspace{1mm}
      With $\alpha=0.01$, $\lambda=1$, $m=50$:
      \[ 1 - \alpha\frac{\lambda}{m} = 1 - \frac{0.01}{50} = 0.9998 \]
      Tiny per step --- but $0.9998^{1000} = 0.8187$. Left alone for a thousand iterations, a
      weight decays to $82\%$ of its value.
    \end{column}
    \begin{column}{0.44\textwidth}
      \begin{keyidea}
        \footnotesize
        This is why L2 regularization is also called \textbf{weight decay}. Each weight leaks
        away a little every iteration, and survives only if the data keeps pushing it back
        up. A feature that does not help is quietly driven towards zero.
      \end{keyidea}
      \par\vspace{1mm}
      The shrink factor must stay in $(0,1)$: if $\alpha\lambda/m > 1$ the sign of $w_j$ flips
      each iteration and training diverges.
    \end{column}
  \end{columns}
\end{frame}

% ===============================================================
\section{Regularized Logistic Regression}

\begin{frame}{The same penalty, the same update}
  \small
  Add the identical term to the logistic cost of Part 1:
  \[
    J(\vw,b) = -\frac{1}{m}\sum_{i=1}^{m}
      \Bigl[\yii{i}\log f^{(i)} + (1-\yii{i})\log\bigl(1-f^{(i)}\bigr)\Bigr]
      \;+\; \frac{\lambda}{2m}\sum_{j=1}^{n}w_j^2 ,
  \]
  with $f^{(i)} = g(\vw\cdot\xii{i}+b)$. Differentiating gives
  \[
    \frac{\partial J}{\partial w_j}
      = \frac{1}{m}\sum_{i=1}^{m}\bigl(f^{(i)}-\yii{i}\bigr)\xjj{i}{j} + \frac{\lambda}{m}w_j ,
    \qquad
    \frac{\partial J}{\partial b}
      = \frac{1}{m}\sum_{i=1}^{m}\bigl(f^{(i)}-\yii{i}\bigr) .
  \]
  \begin{columns}[T,onlytextwidth]
    \begin{column}{0.5\textwidth}
      \begin{keyidea}
        \footnotesize
        Character for character the regularized linear-regression gradient. The only
        difference, again, is the definition of $f$. One implementation of
        \texttt{gradient(X, y, w, b, lam)} serves both models.
      \end{keyidea}
    \end{column}
    \begin{column}{0.46\textwidth}
      \footnotesize
      \begin{block}{In code}
        \texttt{dj\_dw = X.T @ err / m + lam * w / m}\\
        \texttt{dj\_db = err.mean()\ \ \ \ \ \ \ \ \# no lam}
      \end{block}
      Geometrically, shrinking $\vw$ in a logistic model flattens the sigmoid along the
      boundary --- predictions become less confident and the boundary straighter.
    \end{column}
  \end{columns}
\end{frame}

\begin{frame}{Choosing $\lambda$, and a note on L1}
  \footnotesize
  \begin{columns}[T,onlytextwidth]
    \begin{column}{0.5\textwidth}
      \begin{block}{Choosing $\lambda$}
        Search on a multiplicative grid, as with $\alpha$:
        \[ 0,\ 0.01,\ 0.03,\ 0.1,\ 0.3,\ 1,\ 3,\ 10,\ \dots \]
        For each, train to convergence and measure the cost on data held out from training.
        Pick the $\lambda$ with the lowest held-out cost --- \emph{never} the lowest training
        cost, which always chooses $\lambda=0$.
      \end{block}
      \begin{pitfall}
        Regularization penalises $w_j$, whose size depends on the units of feature $j$.
        Without feature scaling the penalty falls unequally on features measured in
        square feet and features measured in rooms. \textbf{Scale first, always.}
      \end{pitfall}
    \end{column}
    \begin{column}{0.46\textwidth}
      \begin{block}{L1 --- a different penalty}
        Replace $\sum_j w_j^2$ with $\sum_j |w_j|$ and the behaviour changes qualitatively:
        weights are driven to \emph{exactly} zero rather than merely made small.
        \begin{itemize}
          \item \textbf{L2 (ridge):} shrinks everything, keeps everything.
          \item \textbf{L1 (lasso):} performs feature selection as a side effect --- useful
                when you need to know which features matter.
        \end{itemize}
        L2 is the default and the only one used in this course.
      \end{block}
      \begin{keyidea}
        \footnotesize
        Regularization does not make a bad model good. It buys back some of the flexibility
        you spent on extra features.
      \end{keyidea}
    \end{column}
  \end{columns}
\end{frame}

% ===============================================================
\section{Wrap-up and Assessment}

\begin{frame}{Summary of Week 3, and of Module 1}
  \footnotesize
  \begin{columns}[T,onlytextwidth]
    \begin{column}{0.5\textwidth}
      \begin{block}{Concepts}
        \begin{itemize}
          \item Underfitting is high bias; overfitting is high variance. Zero training cost
                is a symptom, not an achievement.
          \item The goal is generalization, which $J_{\text{train}}$ cannot measure.
          \item Remedies: more data, fewer features, or regularization.
          \item L2 adds $\frac{\lambda}{2m}\sum_j w_j^2$; $b$ is left out.
          \item The update becomes shrink-then-step --- weight decay.
          \item Linear and logistic regression share one regularized gradient.
        \end{itemize}
      \end{block}
    \end{column}
    \begin{column}{0.46\textwidth}
      \begin{block}{Formulas}
        \[ J = \frac{1}{2m}\sum_i\bigl(f^{(i)}-\yii{i}\bigr)^2 + \frac{\lambda}{2m}\sum_{j=1}^{n}w_j^2 \]
        \[ w_j := \Bigl(1-\alpha\tfrac{\lambda}{m}\Bigr)w_j - \tfrac{\alpha}{m}\sum_i\bigl(f^{(i)}-\yii{i}\bigr)\xjj{i}{j} \]
      \end{block}
      \begin{block}{Next: Module 2}
        \footnotesize
        Neural networks. Every idea from Module 1 --- the model/cost/gradient recipe, feature
        scaling, the learning rate, overfitting, regularization --- carries over unchanged.
        Only the model gets bigger.
      \end{block}
    \end{column}
  \end{columns}
\end{frame}

\begin{frame}{Descriptive questions}
  \footnotesize
  \begin{enumerate}
    \item Define underfitting and overfitting, and give the bias/variance name for each.
          State one symptom that distinguishes them. \hfill\clo{6}
    \item Seven data points are fitted by a degree-$6$ polynomial and $J_{\text{train}}=0$.
          Explain why this happens and why it is not good news. \hfill\clo{6}
    \item Explain the term ``high variance'' by describing what happens if the training set
          is redrawn and the model refitted. \hfill\clo{5}
    \item Why can overfitting not be detected from the training cost? What would you need
          instead? \hfill\clo{5}
    \item List the three remedies for overfitting given in this lecture, and state one
          limitation of each. \hfill\clo{6}
    \item Write the regularized cost for linear regression. Explain the role of $\lambda$ and
          why the penalty is divided by $2m$. \hfill\clo{6}
  \end{enumerate}
\end{frame}

\begin{frame}{Descriptive questions (continued)}
  \footnotesize
  \begin{enumerate}\setcounter{enumi}{6}
    \item Why is $b$ excluded from the regularization term? \hfill\clo{4}
    \item Describe the fitted model when $\lambda=0$ and when $\lambda=10^{10}$, and name the
          failure mode in each case. \hfill\clo{6}
    \item Derive $\partial J/\partial w_j$ for the regularized squared-error cost, and state
          how the update rule changes. \hfill\clo{6}
    \item Show that the regularized update can be written as
          $w_j := (1-\alpha\lambda/m)w_j - \alpha\,\partial J_{\text{fit}}/\partial w_j$, and
          explain the name ``weight decay''. For $\alpha=0.01$, $\lambda=1$, $m=50$, compute
          the shrink factor and its value after $1000$ iterations. \hfill\clo{8}
    \item Write the regularized cost and gradient for logistic regression, and state exactly
          what differs from the linear case. \hfill\clo{6}
    \item Why must features be scaled before regularization is applied? Illustrate with two
          features measured in square feet and in number of rooms. \hfill\clo{6}
    \item Contrast L1 and L2 penalties. When would you prefer L1? \hfill\clo{5}
  \end{enumerate}
\end{frame}

\begin{frame}{Design questions}
  \footnotesize
  \begin{enumerate}
    \item \textbf{Design a $\lambda$ search.} You have $m=800$ examples and $n=60$ features.
          Design a complete procedure for choosing $\lambda$: how you split the data, the grid
          you search, what you measure, and the stopping rule. State what you would report to
          justify the chosen value. \hfill\clo{12}
    \item \textbf{Diagnose from symptoms alone.} For each case, decide underfitting or
          overfitting and give two concrete next actions: (a) $J_{\text{train}}=0.02$,
          $J_{\text{new}}=0.9$; (b) $J_{\text{train}}=0.55$, $J_{\text{new}}=0.58$, both far
          above what a human achieves; (c) $J_{\text{train}}=0.31$, $J_{\text{new}}=0.33$, and
          a human achieves $0.30$. \hfill\clo{12}
    \item \textbf{Design under a data constraint.} A rural clinic can supply only $150$
          labelled records but records $80$ measurements per patient. Propose a modelling
          strategy, state which of the three remedies you rely on and why, and identify the
          risk your strategy does not remove. \hfill\clo{12}
    \item \textbf{Design a regularization unit test.} Write tests that would catch: $b$ being
          penalised, the penalty being added to the cost but not the gradient, and $\lambda$
          being applied before feature scaling. For each, give the input and the expected
          observable difference. \hfill\clo{12}
    \item \textbf{Argue the trade-off to a stakeholder.} A hospital director asks why you
          deployed the model with $94\%$ training accuracy rather than the one with $100\%$.
          Write the explanation you would give, in language that uses no equations.
          \hfill\clo{10}
  \end{enumerate}
\end{frame}

\begin{frame}{Reading}
  \begin{block}{Primary text}
    G.~James, D.~Witten, T.~Hastie and R.~Tibshirani,
    \emph{An Introduction to Statistical Learning}, 2nd ed., Springer, 2021.\\[2pt]
    \hfill{\small$\rightarrow$ \textbf{\S2.2} --- assessing model accuracy, the
    bias--variance trade-off, and why training error falls as flexibility rises.}\\[2pt]
    \hfill{\small$\rightarrow$ \textbf{\S6.2.1--6.2.2} --- ridge regression and the lasso:
    the penalty, the role of the tuning parameter, and why the predictors must be
    standardised.}
  \end{block}
  \begin{block}{Supplementary}
    C.~M.~Bishop, \emph{Pattern Recognition and Machine Learning}, Springer, 2006.\\[2pt]
    \hfill{\small$\rightarrow$ \textbf{\S1.1} and \textbf{\S3.1.4} --- polynomial curve
    fitting with a regularization term, and regularized least squares in matrix form.}
  \end{block}
  \vspace{1mm}
  \begin{center}
    \small\color{cugray}
    Lab (CSE 816, Week 3): gradient descent for logistic regression; logistic regression with
    scikit-learn; overfitting; regularization.\\
    Practice quizzes: \emph{Gradient descent for logistic regression}; \emph{The problem of
    overfitting}.
  \end{center}
\end{frame}

\end{document}
